\documentclass[../../../include/open-logic-section]{subfiles}

\begin{document}

\olfileid{sth}{ordinals}{basic}
\olsection{序数的基本性质}

我们观察到，最前面的几个序数就是自然数。
发展序数理论的主要目的是推广自然数上成立的归纳法原理。
我们将通过一系列基本结果来逐步引出这一原理。


\begin{lem}\ollabel{ordmemberord}
一个序数的每个元素都是序数。
\end{lem}

\begin{proof}
设 $\alpha$ 是一个序数且 $b \in \alpha$。由于 $\alpha$ 是传递的，$b \subseteq \alpha$。
既然 $\in$ 良序 $\alpha$，$\in$ 也良序 $b$。

为证 $b$ 是传递的，假设 $x \in c \in b$。由于 $b \subseteq \alpha$，有 $c \in \alpha$。
再由 $\alpha$ 是传递的，$c \subseteq \alpha$，从而 $x \in \alpha$。
所以 $x, c, b \in \alpha$。但 $\in$ 良序 $\alpha$，根据 \olref[sth][ordinals][wo]{wo:strictorder}，$\in$ 是 $\alpha$ 上的传递关系。
因此，由 $x \in c \in b$，可得 $x \in b$。
全称概括即得 $c \subseteq b$
\end{proof}

\begin{cor}\ollabel{ordissetofsmallerord}
对任意序数~$\alpha$，$\alpha = \Setabs{\beta \in \alpha}{\beta \text{ is an ordinal}}$
\end{cor}

\begin{proof}
由 \olref{ordmemberord} 立即得到。
\end{proof}


接下来的两个主要结果
\olref{ordinductionschema} 和 \olref{ordtrichotomy} 的大意是：
序数本身被属于关系良序：


\begin{thm}[超限归纳]\ollabel{ordinductionschema}
对于任意公式 $\phi(x)$：
\[
	\text{若 }\exists \alpha \phi(\alpha)\text{，则 }\exists \alpha(\phi(\alpha)
	\land  (\forall \beta \in \alpha) \lnot \phi(\beta))
\]
其中列出的量词隐含地限制于序数。
\end{thm}

\begin{proof}
%\olref{ordinductionschema1}.
假设对某个序数 $\alpha$ 有 $\phi(\alpha)$。如果 $ (\forall \beta
\in \alpha) \lnot \phi(\beta)$，则证毕。否则，由于
$\alpha$ 是序数，它有一个满足 $\phi$ 的 $\in$-最小元素，而根据 \olref{ordmemberord}，这个元素是一个序数。
%	\olref{ordinductionschema2}. Suppose there is some ordinal
%	$\gamma$ such that $\lnot\phi(\gamma)$. Then by
%	\olref{ordinductionschema1}, there is an $\in$-minimal ordinal
%	$\alpha$ for which $\lnot\phi(\alpha)$. So $(\forall \beta <
%	\alpha) \phi(\beta)$, rendering the antecedent of the conditional
%	false.
\end{proof}
\noindent
注意，我们可以将 \olref{ordinductionschema} 等价地表达为如下模式：
\[
\text{若 }\forall \alpha((\forall \beta \in \alpha)\phi(\beta) \lif
\phi(\alpha))\text{，则 }\forall \alpha\phi(\alpha)
\]
只需在 \olref{ordinductionschema} 中取 $\lnot\phi(\alpha)$，然后进行初等逻辑变形即可。


\begin{thm}[三分法]\ollabel{ordtrichotomy}
对任意序数 $\alpha$ 和 $\beta$，$\alpha \in \beta \lor \alpha = \beta \lor \beta \in \alpha$。
\end{thm}

\begin{proof}
证明采用双重归纳，即两次使用
\olref{ordinductionschema}。
称 $x$ 与 $y$ \emph{可比较}，当且仅当 $x \in y \lor x = y \lor y \in x$。

作归纳假设：$\alpha$ 中的每个序数都与\emph{所有}序数可比较。
进一步假设：$\alpha$ 与 $\beta$ 中的每个序数可比较。
我们将证明 $\alpha$ 与 $\beta$ 可比较。
通过对 $\beta$ 归纳，将推出 $\alpha$ 与所有序数可比较；
从而再通过对 $\alpha$ 归纳，\emph{每个}序数都与\emph{所有}序数可比较，得证。
我们只需假设 $\alpha \notin \beta$ 且 $\beta \notin \alpha$，然后证明 $\alpha = \beta$。

为证 $\alpha \subseteq \beta$，取定 $\gamma \in \alpha$；
根据 \olref{ordmemberord}，这是一个序数。
于是由第一个归纳假设，$\gamma$ 与 $\beta$ 可比较。
但如果 $\gamma = \beta$ 或 $\beta \in \gamma$，则 $\beta \in \alpha$
（必要时利用 $\alpha$ 是传递的这一事实），这与我们的假设矛盾；
所以 $\gamma \in \beta$。从而 $\alpha \subseteq \beta$。

完全类似的推理，利用第二个归纳假设，
可证 $\beta \subseteq \alpha$。
因此 $\alpha = \beta$。
\end{proof}\noindent 这样一来，我们有时会记 $\alpha <\beta$
而非 $\alpha \in \beta$，因为 $\in$ 起着序关系的作用。
这样做没有深层次的原因，仅仅是出于习惯，
并且因为写 $\alpha \leq \beta$ 比写 $\alpha \in
\beta \lor \alpha = \beta$ 更容易。\footnote{我们也可以写 $\alpha
\mathrel{\underline{\in}} \beta$；但那完全是非标准的。}

以下是前面两个结论的两个直接推论，第一个用到了我们的新记号：


\begin{cor}\ollabel{ordordered}
若 $\exists \alpha\phi(\alpha)$，则 $\exists \alpha(\phi(\alpha)
\land \forall \beta(\phi(\beta) \lif \alpha \leq \beta))$。
此外，对任意序数 $\alpha, \beta, \gamma$，既有 $\alpha
\notin \alpha$，又有 $\alpha \in \beta \in \gamma \lif \alpha \in
\gamma$。
\end{cor}

\begin{proof}
完全类似于 \olref[wo]{wo:strictorder}。
\end{proof}

\begin{prob}
补全 \olref[sth][ordinals][basic]{ordtrichotomy} 证明中的“完全类似的推理”部分。
\end{prob}

\begin{cor}\ollabel{corordtransitiveord}
$A$ 是序数当且仅当 $A$ 是一个由序数组成的传递集。
\end{cor}

\begin{proof}
\emph{从左到右。}根据 \olref{ordmemberord}。\emph{从右到左。}
如果 $A$ 是一个由序数组成的传递集，那么由 \olref{ordinductionschema} 和 \olref{ordtrichotomy} 可知，$\in$ 良序 $A$。
\end{proof}

现在，我们将 \olref{ordinductionschema} 和 \olref{ordtrichotomy} 理解为表明 $\in$ 良序了所有序数。然而，由于以下结果，我们必须对这种说法\emph{非常谨慎}：

\begin{thm}[布拉利-福尔蒂悖论]\ollabel{buraliforti}
不存在由所有序数构成的集合
\end{thm}

\begin{proof}
为归谬起见，假设 $O$ 是所有序数的集合。如果 $\alpha \in \beta \in O$，那么根据 \olref{ordmemberord}，$\alpha$ 是一个序数，所以 $\alpha \in O$。因此 $O$ 是传递的，从而由 \olref{corordtransitiveord} 知 $O$ 是一个序数。于是 $O \in O$，这与 \olref{ordordered} 矛盾。
\end{proof}
\noindent
这个结果以 \citeauthor{Burali-Forti1897} 命名。但是，Cantor 在 1899 年——在致 Dedekind的一封信中——首先清楚地看到了假设存在一个由所有序数组成的集合所蕴含的\emph{矛盾}。正如 van Heijenoort 所解释的：
\begin{quote}
  布拉利-福尔蒂本人认为这个矛盾通过\emph{归谬法}确立了以下结果：序数的自然序只是一个偏序。
  \citep[p.~105]{Heijenoort1967}
\end{quote}
撇开布拉利-福尔蒂的错误不谈，我们可以将上述内容总结如下。序数是这样的集合：它们各自被属于关系良序，并且所有序数整体上也被属于关系良序（尽管它们整体并不构成一个集合）。

最后，这里还有一些从 \olref{ordinductionschema} 和 \olref{ordtrichotomy} 推导出的关于序数的基本性质。

\begin{prop}
任何严格递减的序数序列都是有限的。
\end{prop}

\begin{proof}
任何无限严格递减的序数序列 $\alpha_0 > \alpha_1 > \alpha_2 > \ldots$ 都没有 $<$-极小元，这与 \olref{ordinductionschema} 矛盾。
\end{proof}

\begin{prop}\ollabel{ordinalsaresubsets}
对任意序数 $\alpha, \beta$，有 $\alpha \subseteq \beta \lor \beta \subseteq \alpha$。
\end{prop}

\begin{proof}
若 $\alpha \in \beta$，因 $\beta$ 是传递的，故 $\alpha \subseteq \beta$。类似地，若 $\beta \in \alpha$，则 $\beta \subseteq \alpha$。而若 $\alpha = \beta$，那么 $\alpha \subseteq \beta$ 且 $\beta \subseteq \alpha$。因此根据 \olref{ordtrichotomy} 即得证。
\end{proof}

\begin{prop}\ollabel{ordisoidentity}
对任意序数 $\alpha, \beta$，$\alpha = \beta$ 当且仅当 $\ordeq{\alpha}{\beta}$。
\end{prop}

\begin{proof}
序数是良序集；所以这由三分法（\olref[basic]{ordtrichotomy}）和 \olref[iso]{wellordnotinitial} 直接可得。
\end{proof}

\begin{prob}\ollabel{probunionordinalsordinal}
证明：如果 $X$ 的每个成员都是序数，那么 $\bigcup X$ 也是一个序数。
\end{prob}

\end{document}